\pagenumbering{arabic}
\setcounter{page}{1}
\fancyfoot[C]{}
\part{Tổng quan Bài toán}
\label{part:ch1}

Chương này phân tích bài toán nghiệp vụ và các yêu cầu đặt ra cho hệ thống phân tích bán lẻ \mintinline{text}{bigdata_inventory}. Mục~\ref{sec:ch1-context} trình bày bối cảnh ngành bán lẻ giày dép, phát biểu bài toán và các câu hỏi nghiệp vụ cốt lõi. Mục~\ref{sec:users} xác định đối tượng người dùng và các quyết định cần hỗ trợ. Mục~\ref{sec:requirements} liệt kê yêu cầu chức năng và yêu cầu phi chức năng gắn với kiến trúc hệ thống. Mục~\ref{sec:kpi} định nghĩa các chỉ số đánh giá (KPI) và ngưỡng phân loại thuật toán, Mục~\ref{sec:data-overview} mô tả hai tập dữ liệu nguồn, và Mục~\ref{sec:scope} nêu phạm vi cùng các giả định thực hiện.

\section{Bối cảnh và bài toán}
\label{sec:ch1-context}

\subsection{Bối cảnh}
\label{sec:context}
Trong chuỗi bán lẻ đa cửa hàng thuộc nhóm hàng giày dép và trang phục (như \textit{Scholar Footwear}, \textit{Femme Footwear} và \textit{Junior Apparel}), hoạt động kinh doanh phát sinh khối lượng lớn giao dịch bán hàng và biến động tồn kho hằng ngày trên nhiều cửa hàng và danh mục sản phẩm~\cite{smith2026retail}. Dữ liệu lịch sử cho thấy doanh số phân tán qua $40$ cửa hàng và $27$ danh mục sản phẩm, đi kèm các yếu tố đặc thù như biến động theo mùa, tỷ lệ hoàn trả hàng và chênh lệch tốc độ tiêu thụ giữa các mã sản phẩm.

\medskip

Nếu vận hành theo phương thức báo cáo tĩnh hoặc xử lý thủ công trên bảng tính rời rạc, bộ phận quản lý gặp ba khó khăn chính:
\begin{itemize}[topsep=0pt, itemsep=1pt]
	\item \textbf{Độ trễ giám sát bán hàng:} Dữ liệu giao dịch mới không được cập nhật liên tục lên bảng điều khiển, làm chậm khả năng phát hiện xu hướng doanh thu, kênh bán hàng chủ lực và nhóm sản phẩm có tỷ lệ hoàn trả cao.
	\item \textbf{Sai lệch trong dự báo nhu cầu:} Dự báo chỉ dựa vào trung bình giản đơn bỏ qua tính mùa vụ theo tháng và trọng số của các kỳ gần nhất, dẫn đến kế hoạch nhập hàng thiếu chính xác khi nhu cầu tăng hoặc giảm đột ngột~\cite{hyndman2021forecasting}.
	\item \textbf{Mất cân bằng tồn kho:} Khi thiếu cơ chế đối chiếu tự động giữa lượng hàng tồn hiện tại và tốc độ bán gần nhất, chuỗi cửa hàng dễ đồng thời gặp hai rủi ro: đứt hàng (\textit{stockout}) gây mất doanh thu ở nhóm bán chạy, hoặc tồn kho quá mức (\textit{overstock}) gây ứ đọng vốn ở nhóm bán chậm.
\end{itemize}

\subsection{Phát biểu bài toán}
\label{sec:problem}
Bài toán của nhóm là \textbf{xây dựng hệ thống phân tích dữ liệu bán lẻ thời gian thực, dự báo nhu cầu và giám sát tồn kho thông minh (\mintinline{text}{bigdata_inventory})}. Hệ thống hỗ trợ kiến trúc dữ liệu kép (\mintinline{javascript}{DATA_BACKEND=memory|lakehouse|shadow}) đọc $410.506$ bản ghi từ tập tin CSV vào RAM ($2{,}5\text{ s}$) hoặc nạp trực tiếp gói dữ liệu Gold tiền tổng hợp từ đường ống \textit{Medallion Lakehouse} ($16\text{ ms}$), và cung cấp bảng điều khiển trực quan Apache ECharts 6 gồm ba thẻ chức năng: \textbf{EDA \& Streaming}, \textbf{Dự báo Nhu cầu} và \textbf{Tồn kho Thông minh} (Hình~\ref{fig:system_overview}).

\begin{figure}[H]
	\centering
	\mermaidfile[width=0.95\linewidth, keepaspectratio]{diagrams/system_overview.mmd}
	\caption{Tổng quan luồng dữ liệu hai chế độ từ nguồn CSV/Lakehouse đến ba phân hệ của \texttt{bigdata\_inventory}}
	\label{fig:system_overview}
\end{figure}

\medskip

Hệ thống hướng tới ba mục tiêu cụ thể:
\begin{itemize}[topsep=0pt, itemsep=1pt]
	\item \textbf{Giám sát luồng giao dịch SSE và khám phá dữ liệu đa chiều (OLAP Cube):} Đẩy luồng giao dịch thời gian thực qua Server-Sent Events (\mintinline{text}{/events}) có dự phòng polling tự động, kết hợp bộ lọc cắt lớp đa chiều theo Cửa hàng, Nhóm hàng và khoảng tháng trên $7.517$ ô khối lập phương cho bốn biểu đồ ECharts 6.
	\item \textbf{Đối chiếu ba phương pháp dự báo nhu cầu kèm khoảng tin cậy:} Cho phép người dùng chọn cặp danh mục và cửa hàng cùng số kỳ tương lai ($4$--$24$ kỳ) để so sánh \textit{Normal Forecast} (SMA), \textit{Super Forecast} (WMA + mùa vụ) và mô hình học máy \textbf{Hybrid\_GBDT\_Seasonal\_Ensemble} kèm dải khoảng tin cậy $80\%$.
	\item \textbf{Phân loại tồn kho tích lũy, cảnh báo ưu tiên và mô phỏng rủi ro ML:} Đánh giá sức khỏe tồn kho theo số ngày bao phủ nhu cầu $30$ ngày, cộng dồn trạng thái qua SSE/polling, lọc top $10$ cảnh báo ưu tiên và tích hợp bộ mô phỏng xác suất rủi ro \textbf{SMOTENC\_HistGBDT\_Tuned} theo thời gian thực.
\end{itemize}

\medskip

Bảng~\ref{tab:business-questions} liệt kê sáu câu hỏi nghiệp vụ (Q1--Q6) mà hệ thống giải đáp thông qua ba phân hệ giao diện.

\begin{table}[H]
	\centering
	\footnotesize
	\begin{tabular}{c p{11.5cm} p{3.5cm}}
		\toprule
		\textbf{Mã} & \textbf{Câu hỏi nghiệp vụ}                                                                                                                                                               & \textbf{Thẻ hỗ trợ} \\
		\midrule
		Q1          & Các giao dịch bán hàng mới phát sinh theo luồng SSE/polling có thông số (ngày, nhóm hàng, danh mục, cửa hàng, số lượng, doanh thu) như thế nào?                                          & EDA \& Streaming    \\
		Q2          & Xu hướng doanh thu tháng, top $10$ danh mục, cơ cấu kênh và tỷ lệ hoàn trả thay đổi ra sao khi lọc theo Cửa hàng, Nhóm hàng và khoảng tháng?                                             & EDA \& Streaming    \\
		Q3          & Với một danh mục và cửa hàng trong $4$--$24$ kỳ tới, dự báo SMA, WMA+Mùa vụ và mô hình học máy \textit{Hybrid GBDT} (kèm khoảng tin cậy $80\%$) khác biệt thế nào?                       & Dự báo Nhu cầu      \\
		Q4          & Ở những kỳ nào SMA và WMA+Mùa vụ cho kết quả chênh lệch đáng kể (tối đa $8$ kỳ đầu tiên) và thuật toán nào cho mức dự báo cao hơn?                                                       & Dự báo Nhu cầu      \\
		Q5          & Tỷ trọng các mặt hàng \textit{Healthy}, \textit{Low Stock}, \textit{Overstock} trên luồng tích lũy là bao nhiêu và cặp danh mục/cửa hàng nào cần xử lý khẩn cấp?                         & Tồn kho Thông minh  \\
		Q6          & Với các thông số tồn kho giả lập (\mintinline{text}{Qty}, doanh số $30$ ngày, trạng thái giá, đơn giá), xác suất rủi ro mô hình \textit{SMOTENC} và cảnh báo hàng tồn chết là bao nhiêu? & Tồn kho Thông minh  \\
		\bottomrule
	\end{tabular}
	\caption{Câu hỏi nghiệp vụ trọng tâm của hệ thống}
	\label{tab:business-questions}
\end{table}

\section{Người dùng và Quyết định cần hỗ trợ}
\label{sec:users}
Hình~\ref{fig:decision_levels} tóm tắt ba nhóm người dùng, các cấp quyết định tương ứng và các thành phần giao diện phục vụ từng cấp. Hệ thống phục vụ quyết định quản trị bán lẻ mang tính bán cấu trúc (\textit{semi-structured decision})~\cite{turban2010dsbi}: máy chủ tự động hóa khâu tổng hợp dữ liệu lớn, tính toán dự báo và phân loại trạng thái, còn người quản lý dựa trên kết quả trực quan để đưa ra quyết định điều phối nhập hàng hoặc xả tồn. Bảng~\ref{tab:user-decision-levels} trình bày chi tiết từng nhóm người dùng và thành phần hỗ trợ trên ứng dụng web.

\begin{figure}[H]
	\centering
	\mermaidfile[width=0.95\linewidth, keepaspectratio]{diagrams/decision_levels.mmd}
	\caption{Các cấp quyết định và thành phần giao diện hỗ trợ người dùng}
	\label{fig:decision_levels}
\end{figure}

\begin{table}[H]
	\centering
	\footnotesize
	\begin{tabular}{p{2.6cm} p{1.9cm} p{4.6cm} p{5.1cm}}
		\toprule
		\textbf{Người dùng}                  & \textbf{Cấp quyết định} & \textbf{Quyết định cần hỗ trợ}                                                                      & \textbf{Thành phần hỗ trợ trên web}                                                                                                          \\
		\midrule
		Ban điều hành / Giám đốc kinh doanh  & Chiến lược, giám sát    & Đánh giá hiệu quả tổng thể, kênh phân phối và kiểm soát tỷ lệ hàng hoàn trả                         & Bốn biểu đồ EDA trên thẻ \textbf{EDA \& Streaming}; biểu đồ tròn tỷ lệ tồn kho tích lũy trên thẻ \textbf{Tồn kho Thông minh}                 \\
		\midrule
		Chuyên viên lập kế hoạch cung ứng    & Chiến thuật             & Xác định sản lượng nhập bổ sung theo từng danh mục và cửa hàng trong $4$--$24$ kỳ                   & Biểu đồ đường kép và bảng đối soát $8$ kỳ sai khác trên thẻ \textbf{Dự báo Nhu cầu}                                                          \\
		\midrule
		Quản lý cửa hàng / Điều phối tồn kho & Tác nghiệp              & Theo dõi luồng bán hàng tức thời; đặt hàng khẩn cấp cho mã thiếu hụt hoặc khuyến mãi xả hàng ứ đọng & Bảng giao dịch trực tiếp ($100$ dòng) trên \textbf{EDA \& Streaming}; $3$ thẻ KPI và bảng Top $10$ cảnh báo trên \textbf{Tồn kho Thông minh} \\
		\bottomrule
	\end{tabular}
	\caption{Người dùng, cấp quyết định và thành phần hệ thống hỗ trợ}
	\label{tab:user-decision-levels}
\end{table}

\section{Yêu cầu nghiệp vụ}
\label{sec:requirements}

\subsection{Yêu cầu chức năng}
\label{sec:functional-requirements}
Bảng~\ref{tab:functional-requirements} tổng hợp $12$ yêu cầu chức năng (FR-01 đến FR-12) được hiện thực xuyên suốt từ đường ống Medallion Lakehouse, tầng phục vụ Node.js/ML đến giao diện React 18 + ECharts 6, gắn với các câu hỏi nghiệp vụ Q1--Q6.

\begin{table}[H]
	\centering
	\footnotesize
	\setlength{\tabcolsep}{3.5pt}
	\begin{tabular}{c p{9.1cm} p{3.7cm} c}
		\toprule
		\textbf{Mã} & \textbf{Mô tả yêu cầu}                                                                                                                                                                                                                                           & \textbf{Dịch vụ / Giao diện}                                                        & \textbf{Câu hỏi} \\
		\midrule
		FR-01       & Hỗ trợ khởi tạo hai chế độ: nạp luồng $410.506$ dòng CSV vào RAM (nhường sự kiện mỗi $30.000$ dòng) hoặc nạp gói Gold \mintinline{text}{serving_bundle.json} của Lakehouse ($< 100\text{ ms}$).                                                                  & \mintinline{javascript}{dataRepository}, \mintinline{javascript}{dataLoaderService} & Q1--Q6           \\
		FR-02       & Cung cấp API kiểm tra trạng thái sẵn sàng (\mintinline{text}{GET /api/health} kèm tên chế độ backend) và chặn truy cập khi chưa sẵn sàng bằng mã \mintinline{text}{HTTP 503}.                                                                                    & \mintinline{javascript}{healthController}, \mintinline{javascript}{readinessGuard}  & --               \\
		FR-03       & Đẩy luồng giao dịch qua SSE (\mintinline{text}{/api/stream/sales/events}, $20$ dòng/$2\text{ s}$, nhịp tim $15\text{ s}$) kèm dự phòng HTTP polling khi có $\ge 2$ lỗi, giữ tối đa $100$ dòng.                                                                   & \mintinline{javascript}{salesStreamService}, \textbf{EdaTab}                        & Q1               \\
		FR-04       & Phục vụ $4$ biểu đồ EDA ECharts 6 hỗ trợ cắt lớp OLAP Cube (\mintinline{text}{GET /api/eda/filters} và tham số \mintinline{text}{store}, \mintinline{text}{division}, \mintinline{text}{startMonth}, \mintinline{text}{endMonth}).                               & \mintinline{javascript}{edaService}, \textbf{EdaTab}                                & Q2               \\
		FR-05       & Cung cấp danh sách danh mục và cửa hàng duy nhất, sắp xếp thứ tự qua \mintinline{text}{/api/forecast/*} (duy trì bí danh tương thích \mintinline{text}{/api/categories}, \mintinline{javascript}{/api/stores}).                                                  & \mintinline{javascript}{forecastController}, \textbf{ForecastTab}                   & Q3               \\
		FR-06       & Tính song song SMA-4, WMA+Mùa vụ ($4$--$24$ kỳ, LRU $1.500$ mục) và chuỗi dự báo học máy \textit{Hybrid GBDT} kèm khoảng tin cậy $80\%$ (\mintinline{text}{mlLower80}, \mintinline{text}{mlUpper80}).                                                            & \mintinline{javascript}{forecastService}, \mintinline{javascript}{mlService}        & Q3               \\
		FR-07       & So sánh hai chuỗi SMA và WMA+Mùa vụ sau chuẩn hóa \mintinline{javascript}{toFixed(2)}, trích xuất tối đa $8$ kỳ sai khác đầu tiên kèm độ lệch và thuật toán cao hơn.                                                                                             & \mintinline{javascript}{forecastService}, \textbf{ForecastTab}                      & Q4               \\
		FR-08       & Phát luồng tồn kho lô $50$ bản ghi mỗi $2\text{ s}$ qua SSE/polling (\mintinline{text}{/api/inventory/stream}), phân loại $3$ mức và cộng dồn nguyên tử vào RAM (kèm \mintinline{javascript}{POST /api/inventory/reset}).                                        & \mintinline{javascript}{inventoryService}, \textbf{InventoryTab}                    & Q5               \\
		FR-09       & Lọc danh sách Top $10$ cảnh báo tồn kho đã khử trùng lặp cặp \textit{(Category, Store)} (\mintinline{text}{GET /api/inventory/alerts}), ưu tiên \textit{Low Stock} trước \textit{Overstock}.                                                                     & \mintinline{javascript}{inventoryService}, \textbf{InventoryTab}                    & Q5               \\
		FR-10       & Vận hành đường ống \textit{Medallion Lakehouse} (\textit{Bronze--Silver--Gold} trên DuckDB/Parquet), kiểm định Schema/DLQ, $8/8$ cổng chất lượng và kiểm toán PSI/KS drift (\mintinline{text}{pipelines/}).                                                      & \mintinline{text}{run_pipeline.py}, \mintinline{text}{dagster_defs}                 & Q1--Q6           \\
		FR-11       & Huấn luyện ngoại tuyến \textit{Hybrid GBDT} + \textit{SMOTENC}/\textit{ADASYN} (\mintinline{text}{ml/}) và phục vụ suy luận tức thời qua \mintinline{text}{/api/ml/metrics}, \mintinline{text}{/api/ml/forecast}, \mintinline{text}{/api/ml/classify-inventory}. & \mintinline{javascript}{mlService}, \textbf{InventoryTab}                           & Q3, Q6           \\
		FR-12       & Trang bị giao diện Sáng/Tối chống chớp khung hình, đồng bộ tham số \mintinline{text}{?tab=...} trên URL, bảng lệnh nhanh \texttt{Cmd+K}/\texttt{Ctrl+K} và khung xương tải dữ liệu (\textit{Skeleton}).                                                          & \mintinline{javascript}{AppContext}, \mintinline{javascript}{CommandPalette}        & Q1--Q6           \\
		\bottomrule
	\end{tabular}
	\caption{Danh sách $12$ yêu cầu chức năng của hệ thống}
	\label{tab:functional-requirements}
\end{table}

\subsection{Yêu cầu phi chức năng}
\label{sec:nfr}
Bên cạnh các chức năng nghiệp vụ, hệ thống được thiết kế để đáp ứng năm nhóm yêu cầu phi chức năng trọng yếu:
\begin{itemize}[topsep=0pt, itemsep=1pt]
	\item \textbf{Hiệu năng khởi tạo và không nghẽn luồng sự kiện:} Ở chế độ \mintinline{text}{memory}, bộ phân tích luồng CSV chủ động nhường vòng lặp sự kiện (\mintinline{javascript}{setImmediate}) sau mỗi $30.000$ bản ghi giúp API \mintinline{text}{/api/health} phản hồi tức thì ($2{,}5\text{ s}$ nạp xong); ở chế độ \mintinline{javascript}{lakehouse}, \mintinline{javascript}{LakehouseRepository} nạp gói Gold tiền tổng hợp trong $16\text{ ms}$ (nhanh gấp $157\times$).
	\item \textbf{Truyền dữ liệu thời gian thực bền vững (SSE \& Fallback):} Giao tiếp thời gian thực ưu tiên \mintinline{javascript}{EventSource} (Server-Sent Events kèm nhịp tim $15\text{ s}$ và cấu hình Nginx \mintinline{nginx}{proxy_buffering off}, \mintinline{nginx}{proxy_read_timeout 3600s}); khi gặp $\ge 2$ lỗi kết nối liên tiếp hoặc ẩn thẻ trình duyệt (\mintinline{javascript}{document.hidden}), hook tự động chuyển về HTTP polling $2.000\text{ ms}$ có \mintinline{javascript}{AbortController} ($10.000\text{ ms}$ timeout).
	\item \textbf{Độ sẵn sàng và bảo vệ truy cập:} Lớp trung gian \mintinline{javascript}{readinessGuard} chặn mọi đường dẫn \mintinline{text}{/api/*} (ngoại trừ \mintinline{text}{/api/health}) với mã HTTP \mintinline{text}{503 SERVICE_UNAVAILABLE} trong lúc khởi tạo, bảo vệ API tránh đọc trạng thái chưa hoàn chỉnh.
	\item \textbf{Tính mô-đun và khả năng bảo trì:} Phần máy chủ tách biệt theo \textit{Clean Architecture} và mẫu \textit{Repository} (\mintinline{text}{backend/repositories/}) giữa logic tính toán thuần (\mintinline{text}{algorithms/}), tầng dịch vụ (bao gồm \mintinline{javascript}{mlService}) và tầng HTTP. Giao diện tách ngữ cảnh \mintinline{javascript}{AppContext} (quản lý thẻ, trạng thái kết nối, Sáng/Tối và đồng bộ URL) khỏi \mintinline{javascript}{ToastContext}.
	\item \textbf{Đóng gói đa container và khả năng mở rộng:} \mintinline{text}{docker-compose.yml} đóng gói container \mintinline{text}{backend} (Node.js cổng $3000$ kèm \mintinline{text}{HEALTHCHECK} $5\text{ s}$ và sẵn sàng \mintinline{text}{ml/artifacts/}, \mintinline{javascript}{data/lakehouse/}) cùng container \mintinline{text}{frontend} (Nginx Alpine cổng $80 \rightarrow 3001$); đồng thời cung cấp hồ sơ \mintinline{bash}{--profile pipeline} tích hợp Redpanda, Redis $7$ và ClickHouse khi cần thử nghiệm hạ tầng streaming.
\end{itemize}

\section{Chỉ số đánh giá (KPI)}
\label{sec:kpi}
Các chỉ số và ngưỡng quyết định của hệ thống được định nghĩa tập trung trong \mintinline{text}{constants.js} (thư mục \mintinline{text}{backend/config/}), \mintinline{javascript}{dataLoaderService.js} (thư mục \mintinline{text}{backend/services/}) và ba hàm thuật toán thuần túy tại \mintinline{text}{backend/algorithms/}. Nhờ đồng bộ hằng số giữa máy chủ và giao diện, số liệu hiển thị trên ba thẻ của ứng dụng khớp hoàn toàn với kết quả tính toán ở tầng dịch vụ.

\subsection{Quy ước Phạm vi tính toán}
\label{sec:kpi-scope}

\begin{figure}[H]
	\centering
	\mermaidfile[width=0.95\linewidth, keepaspectratio]{diagrams/kpi_scope.mmd}
	\caption{Phạm vi tính toán chỉ số và ngưỡng phân loại tồn kho}
	\label{fig:kpi_scope}
\end{figure}

Tùy theo đặc thù bài toán, các chỉ số được tính trên ba phạm vi sau:
\begin{itemize}
	\item \textbf{Toàn bộ lịch sử giao dịch:} Tính trên $125.751$ bản ghi bán hàng (từ \texttt{2025-06-01} đến \texttt{2026-04-24}) phục vụ bốn biểu đồ EDA và xây dựng chuỗi thời gian theo tháng cho từng cặp danh mục -- cửa hàng.
	\item \textbf{Cửa sổ $30$ ngày gần nhất (\mintinline{text}{SALES_LOOKBACK_DAYS = 30}):} Lấy mốc ngày giao dịch lớn nhất trong tập dữ liệu (\texttt{2026-04-24}) lùi về $30$ ngày để tổng hợp sản lượng bán $30$ ngày (\mintinline{text}{prodSales30d}) của từng mã sản phẩm, làm căn cứ tính tốc độ tiêu thụ bình quân ngày.
	\item \textbf{Luồng tích lũy trên bộ nhớ (\textit{Cumulative Stream}):} Cộng dồn trạng thái phân loại qua các lô ngẫu nhiên ($20$ dòng/lần cho bảng bán hàng; $50$ bản ghi/lần cho biểu đồ tồn kho), tồn tại xuyên suốt vòng đời máy chủ cho đến khi người dùng đặt lại.
\end{itemize}

\subsection{Định nghĩa Chỉ số và Ngưỡng thuật toán}
\label{sec:kpi-main}

\begin{table}[H]
	\centering
	\footnotesize
	\begin{tabular}{c p{2.5cm} p{8cm} p{3.3cm}}
		\toprule
		\textbf{Phân hệ} & \centering\textbf{Chỉ số / Thuật toán}             & \centering\textbf{Công thức và Ngưỡng xác định}                                                                                                                                                          & \textbf{Hiển thị trên web}                   \\
		\midrule
		EDA              & Doanh thu tháng                                    & Tổng \mintinline{text}{Sales Amount} nhóm theo định dạng \mintinline{text}{YYYY-MM}, sắp xếp tăng dần theo thời gian                                                                                     & Xu hướng doanh thu (\$)                      \\
		EDA              & Top $10$ danh mục                                  & Tổng \mintinline{text}{Qty Sold} theo \mintinline{text}{Product Category}, sắp xếp giảm dần và lấy $10$ danh mục dẫn đầu                                                                                 & Top $10$ danh mục bán chạy                   \\
		EDA              & Cơ cấu kênh bán hàng                               & Tổng \mintinline{text}{Sales Amount} theo \mintinline{text}{Sales Channel} và tỷ trọng phần trăm trên tổng doanh thu                                                                                     & Phân bổ doanh thu theo kênh                  \\
		EDA              & Tỷ lệ hoàn trả                                     & Số giao dịch có \mintinline{text}{Is Return = 1} chia tổng giao dịch của mỗi \mintinline{text}{Product Division} ($\times 100\%$, làm tròn $2$ chữ số)                                                   & Tỷ lệ đổi trả (\%)                           \\
		\midrule
		Dự báo           & Normal Forecast (SMA)                              & Trung bình cộng trượt của $N=4$ kỳ gần nhất: $\hat{y}_{t} = \dfrac{1}{4}\displaystyle\sum_{k=1}^{4} y_{t-k}$, làm tròn $2$ chữ số thập phân                                                              & Đường xanh (\textit{Normal})                 \\
		Dự báo           & Super Forecast (WMA + Mùa vụ)                      & WMA trọng số $w=[1, 2, 3, 4]$ trên $4$ kỳ gần nhất, nhân hệ số mùa vụ $SF_m = \overline{y}_m / \overline{y}$ của tháng dự báo $m$, làm tròn $2$ chữ số                                                   & Đường cam (\textit{Super})                   \\
		Dự báo           & Chênh lệch đối soát                                & $\Delta = |\hat{y}_{\text{super}} - \hat{y}_{\text{normal}}|$ và $\%\Delta = \dfrac{\Delta}{|\hat{y}_{\text{normal}}|} \times 100\%$; lọc tối đa $8$ kỳ có \mintinline{javascript}{toFixed(2)} khác nhau & Bảng đối soát $8$ kỳ                         \\
		\midrule
		Tồn kho          & Tốc độ bán ngày (\mintinline{text}{avgDailySales}) & $\text{sales30d} / 30$, với \mintinline{text}{sales30d} là tổng \mintinline{text}{Qty Sold} của sản phẩm trong $30$ ngày gần nhất                                                                        & Tham số nội suy ngưỡng                       \\
		Tồn kho          & \textbf{Healthy}                                   & $\text{avgDailySales} \le 0$ hoặc $\text{qtyOnHand} \ge 14 \times \text{avgDailySales}$ (và $< 60 \times \text{avgDailySales}$)                                                                          & Thẻ xanh lá (\mintinline{text}{Cover > 14d}) \\
		Tồn kho          & \textbf{Low Stock}                                 & $\text{avgDailySales} > 0$ và $\text{qtyOnHand} < 14 \times \text{avgDailySales}$                                                                                                                        & Thẻ cam (\mintinline{text}{Cover < 14d})     \\
		Tồn kho          & \textbf{Overstock}                                 & $\text{avgDailySales} > 0$ và $\text{qtyOnHand} \ge 60 \times \text{avgDailySales}$                                                                                                                      & Thẻ đỏ (\mintinline{text}{Cover > 60d})      \\
		\midrule
		Học máy          & Hybrid GBDT + $80\%$ CI                            & Kết hợp phân vị $q=0{,}42$ (\textit{HistGBDT}) với \textit{Ridge} \& hệ số mùa vụ; dải tin cậy biên phân vị $10\%$--$90\%$ hiệu chuẩn trên tập OOS                                                       & Đường tím \& Dải $80\%$ CI                   \\
		Học máy          & Rủi ro ML (\textit{SMOTENC})                       & Xác suất $P(\text{Healthy/Low/Over})$ từ hàm logistic hiệu chuẩn kèm cờ \mintinline{text}{isDeadStockCandidate} khi $\text{avgDailySales} \le 0$                                                         & Bộ mô phỏng rủi ro (\textbf{Tab 3})          \\
		Chất lượng       & Kiểm soát trôi dạt (\textit{Drift})                & Chỉ số ổn định quần thể ($\text{PSI} < 0{,}2$) và kiểm định hai mẫu Kolmogorov--Smirnov (\textit{KS}) giữa tập tham chiếu ($\le$ \texttt{2026-01}) và tập mới ($\ge$ \texttt{2026-02})                   & Báo cáo Gold Lakehouse                       \\
		\bottomrule
	\end{tabular}
	\caption{Định nghĩa các chỉ số và ngưỡng đánh giá trên ba phân hệ và mô hình ML}
	\label{tab:kpi-definitions}
\end{table}

Trên dữ liệu lịch sử thực tế ($125.751$ giao dịch bán hàng), hệ thống ghi nhận tổng doanh thu đạt $10.482.110{,}25$ đơn vị tiền tệ với tổng số lượng bán $116.995$ đơn vị và $7.491$ lượt hoàn trả (tương ứng tỷ lệ hoàn trả trung bình $5{,}96\%$). Trên $284.755$ bản ghi tồn kho, tổng lượng hàng tồn kho ghi nhận đạt $653.407$ đơn vị.

\section{Tổng quan dữ liệu}
\label{sec:data-overview}
Hệ thống sử dụng hai tập tin CSV từ bộ dữ liệu \textit{Retail Transactions and Stocks Data}~\cite{smith2026retail} đặt trong thư mục \mintinline{text}{data/}, gồm tổng cộng $410.506$ bản ghi. Bảng~\ref{tab:data-overview} tóm tắt quy mô, khoảng thời gian và các trường dữ liệu chính được hệ thống trích xuất.

\begin{table}[H]
	\centering
	\footnotesize
	\begin{tabular}{c c c p{1.6cm} p{5.8cm}}
		\toprule
		\textbf{Tập tin nguồn}                         & \textbf{Số dòng} & \textbf{Số cột} & \centering\textbf{Thời gian}                                     & \textbf{Trường dữ liệu sử dụng chính}                                                                                                                                                                                                                                                                \\
		\midrule
		\mintinline{text}{retail_sales_ml_apl.csv}     & $125.751$        & $17$            & \texttt{2025/06/01} $\rightarrow$ \texttt{26/04/24}              & \mintinline{text}{Transaction Date}, \mintinline{text}{Is Return}, \mintinline{text}{Product No}, \mintinline{text}{Product Division}, \mintinline{text}{Product Category}, \mintinline{text}{Store}, \mintinline{text}{Sales Channel}, \mintinline{text}{Qty Sold}, \mintinline{text}{Sales Amount} \\
		\midrule
		\mintinline{text}{retail_inventory_ml_apl.csv} & $284.755$        & $18$            & \texttt{2025/06/01} $\rightarrow$  \mintinline{text}{9999/12/31} & \mintinline{text}{Start Date}, \mintinline{text}{End Date}, \mintinline{text}{Stock Status}, \mintinline{text}{Product No}, \mintinline{text}{Product Category}, \mintinline{text}{Store}, \mintinline{text}{Store Type}, \mintinline{text}{Qty on hand}                                             \\
		\bottomrule
	\end{tabular}
	\caption{Tổng quan hai tập dữ liệu nguồn trong hệ thống}
	\label{tab:data-overview}
\end{table}

Phân tích cấu trúc hai tập tin nguồn cho thấy năm đặc điểm dữ liệu ảnh hưởng trực tiếp tới thiết kế xử lý của \mintinline{javascript}{dataLoaderService}:
\begin{itemize}[topsep=0pt, itemsep=1pt]
	\item \textbf{Phân cấp sản phẩm và cửa hàng đồng nhất:} Cả hai tập tin đều phủ trên $40$ cửa hàng (\texttt{STR-1006} trở đi) và $27$ danh mục sản phẩm thuộc $3$ nhóm hàng (\textit{Scholar Footwear}, \textit{Femme Footwear}, \textit{Junior Apparel}), liên kết thông qua mã sản phẩm \mintinline{text}{Product No} ($2.326$ mã trong bảng tồn kho) và mã cửa hàng \mintinline{text}{Store}.
	\item \textbf{Dữ liệu giao dịch một kênh và cửa hàng một loại:} Toàn bộ bản ghi bán hàng thuộc một kênh phân phối (\textit{Channel Alpha}) và toàn bộ bản ghi tồn kho thuộc loại hình \textit{Retail Storefront}; điều này cho phép kiểm chứng cấu trúc biểu đồ phân bổ mà không làm phân mảnh tập mẫu.
	\item \textbf{Trạng thái giá và mốc kết thúc hiệu lực:} Tập tồn kho phân biệt $5$ trạng thái định giá (\textit{Full Price}, \textit{Markdown Tier 1}, \textit{Markdown Tier 2}, \textit{Clearance}, \textit{Promo}) và sử dụng mốc giả \mintinline{text}{End Date = 9999-12-31} cho các bản ghi tồn kho còn hiệu lực.
	\item \textbf{Đặc thù dữ liệu hoàn trả:} Cột \mintinline{text}{Reason of Return} bị bỏ trống ở toàn bộ $125.751$ giao dịch ($100\%$), kể cả trên $7.491$ giao dịch có \mintinline{text}{Is Return = 1}; do đó chỉ số hoàn trả được tính hoàn toàn theo tỷ lệ nhị phân trên \mintinline{text}{Is Return}.
	\item \textbf{Chuẩn hóa định dạng số và chuỗi CSV:} Các trường số thực như \mintinline{text}{Qty Sold}, \mintinline{text}{Sales Amount} và \mintinline{text}{Cogs} lưu dưới dạng chuỗi có $4$ chữ số thập phân trong CSV (ví dụ \texttt{115.0000}), đòi hỏi tầng nạp ép kiểu số an toàn và sử dụng bộ phân tách CSV chuẩn RFC 4180 (\mintinline{text}{csv-parser} kèm dự phòng \mintinline{javascript}{parseCSVLine}) để xử lý trường văn bản có dấu ngoặc kép.
\end{itemize}

\section{Phạm vi và Giả định}
\label{sec:scope}

\subsection*{Trong phạm vi}
\label{sec:in-scope}
\begin{itemize}[topsep=0pt, itemsep=1pt]
	\item \textbf{Xử lý hai chế độ (\mintinline{text}{memory} và \mintinline{javascript}{lakehouse}):} Tự động nạp $410.506$ bản ghi CSV vào RAM hoặc đọc gói Gold từ đường ống \textit{Medallion Lakehouse} (\textit{Bronze--Silver--Gold} DuckDB/Parquet), kèm $8/8$ cổng chất lượng và kiểm định trôi dạt PSI/KS.
	\item \textbf{Luồng SSE, cắt lớp OLAP Cube và dự báo ML:} Cung cấp luồng thời gian thực qua SSE (\mintinline{text}{/events}) có dự phòng polling $2\text{ s}$, cắt lớp $7.517$ ô OLAP Cube trên thẻ EDA, đối chiếu SMA, WMA+mùa vụ và \textit{Hybrid GBDT} (+ khoảng tin cậy $80\%$), cùng mô phỏng rủi ro \textit{SMOTENC}.
	\item \textbf{Kiểm thử đa tầng và đóng gói Docker:} Vận hành bộ kiểm thử đơn vị/E2E ($129$ bài), chịu tải $7$+4B pha (\mintinline{javascript}{stress_test.js}), đối đẳng Lakehouse (\mintinline{text}{shadow_parity.js}), kiểm thử Python (\mintinline{bash}{pytest}) và đóng gói React 18 + Apache ECharts 6 qua Docker Compose/Nginx.
\end{itemize}

\subsection*{Ngoài phạm vi}
\label{sec:out-of-scope}
\begin{itemize}[topsep=0pt, itemsep=1pt]
	\item \textbf{Tiến trình tiêu thụ bất đồng bộ trực tuyến từ Message Broker:} Dù tệp cấu hình \mintinline{text}{docker-compose.yml} hỗ trợ hồ sơ \mintinline{bash}{--profile pipeline} khởi chạy Redpanda, Redis và ClickHouse cùng biến môi trường kết nối, máy chủ Node.js hiện tại đọc dữ liệu từ CSV trên RAM hoặc tệp Gold Lakehouse chứ chưa duy trì vòng lặp consumer/sink ghi-đọc trực tiếp trên dây mạng.
	\item \textbf{Ghi ngược xuống dữ liệu nguồn và đặt hàng tự động:} Hệ thống chỉ hỗ trợ đọc, phân tích, dự báo và cảnh báo; không tự động tạo đơn đặt hàng tới nhà cung cấp hay cập nhật sửa đổi tập tin CSV gốc.
	\item \textbf{Xác thực và phân quyền người dùng:} Phiên bản hiện tại phục vụ trình diễn phân tích tập trung nên chưa tích hợp đăng nhập hoặc phân quyền theo vai trò cửa hàng.
\end{itemize}

\subsection*{Giả định}
\label{sec:assumptions}
\begin{enumerate}[label=\textit{\arabic*},topsep=0pt, itemsep=1pt]
	\item \textbf{Cửa sổ doanh số đại diện:} Tốc độ tiêu thụ bình quân ngày của mỗi sản phẩm được ước lượng bằng tổng \mintinline{text}{Qty Sold} trong $30$ ngày gần nhất tính từ ngày giao dịch mới nhất trong tập dữ liệu (\texttt{2026-04-24}) chia cho $30$.
	\item \textbf{Chu kỳ dự báo theo tháng:} Dữ liệu đầu vào của thuật toán dự báo được tổng hợp theo từng tháng (\mintinline{text}{YYYY-MM}) và điền giá trị $0$ cho các tháng khuyết trên toàn bộ $11$ tháng lịch sử; cặp không có lịch sử giao dịch sẽ dùng chuỗi $0$ đủ $11$ tháng.
	\item \textbf{Mặc định với sản phẩm chưa có lịch sử bán:} Nếu một mã sản phẩm trong kho không phát sinh doanh số trong $30$ ngày gần nhất ($\text{avgDailySales} \le 0$), hàm \mintinline{javascript}{classifyInventory} mặc định phân loại là \textbf{Healthy} để tránh báo động giả thiếu hàng.
	\item \textbf{Thứ tự ưu tiên cảnh báo:} Bảng Top $10$ cảnh báo tồn kho khử trùng lặp theo cặp \textit{(Category, Store)}, giữ tối đa $3$ cảnh báo \textit{Overstock} (nếu có) và sắp xếp toàn bộ \textit{Low Stock} lên trước \textit{Overstock}.
	\item \textbf{Trạng thái tích lũy trên máy chủ:} Bộ đếm \mintinline{javascript}{cumulativeState} trên \mintinline{javascript}{InventoryService} giữ nguyên trong RAM kể cả khi tải lại trang trình duyệt, chỉ trở về $0$ khi máy chủ khởi động lại hoặc nhận yêu cầu \mintinline{javascript}{POST /api/inventory/reset}.
\end{enumerate}
